Zusammenfassend werden wir zeigen:

\emph{Satz III:} Mit den Identit\"aten
\begin{equation}
 \pair{q_\rho A^\sigma}{p_\tau B_{\rho+\sigma-\tau}}
 =\sum_{\substack{0,\tau-\sigma}}^{\substack{\tau,\rho}}
 \eps\pair{R^{\tau-\alpha}q_{n-\tau}}{R_{\rho-\alpha}p_{n-\rho}},
\tag{36}
\end{equation}
wo die Reihen $R$ durch (33.) bestimmt sind, und mit den durch Aufl\"osung von $q_\rho,p_\tau$ in $q_{\rho_1}C^{\sigma_1}$ usf. daraus entstehenden ist die Gesamtheit der unzerlegbaren Identit\"aten ersch\"opft.

In der Tat sind die entstehenden Identit\"aten die allgemeinsten ihrer Art, da die einzelnen Reihen in bezug auf Gewicht und Anzahl keiner Beschr\"ankung mehr unterliegen, wie dies noch bei (20.), (21.) der Fall war. Es existieren aber auch keine weiteren Identit\"aten zwischen diesen Reihen; denn bei Aufl\"osung einer Reihe $p_\tau$ in $y^{(1)}\cdots y^{(\tau)}$ sind die allgemeinsten Identit\"aten aus (20.), also nach \S\ 3 (Schlu\ss) aus (20a.) zusammengesetzt; und zwar ist die im Anschlu\ss\ an (16.) diskutierte Zusammensetzung genau die in \S\ 4 durchgef\"uhrte. Den \"Ubergang zu beliebig vielen Reihen liefert aber, wie die Diskussion von (20a.) zeigt, die Anwendung immer derselben Operation auf die durch Aufl\"osung von $q_\rho$ in $q_{\rho_1}C^{\sigma_1}$ usf. entstehenden Ausdr\"ucke. Es folgt daraus auch, da\ss\ der direkte \"Ubergang von (20.) an Stelle von (20a.) nichts Neues liefert; dies l\"a\ss t sich rechnerisch leicht best\"atigen, indem man einmal eine Reihe $q_\rho$ in (25.) spezialisiert, das andere Mal nach der Methode von \S\ 4 den \"Ubergang von (20.) macht; es entstehen beide Male die gleichen Summen, die durch Anwendung des allgemeinen Produktsatzes (siehe oben) in die aus (36.) sich ergebenden Identit\"aten \"ubergehen. Die Identit\"aten (20.), (21.) entsprechen den Spezialf\"allen $\tau=1$, $\rho=1$; es treten dann nur je zwei Einzelsummen auf, also einer obigen Bemerkung entsprechend explizite Darstellung ohne die Substitution (33.), was mit dem fr\"uheren \"ubereinstimmt.
